\scanpage{3}
\deleted{
deux éléments de $L$ (\uline{adjacents} l'un de l'autre) \inserted{appelés ``lieux \unsure{extrémals|extrémales} du segment''}

\textcircled{8} Une partie $\tilde I$ de $L$ (l'ens des lieux de $\mathcal X$. \uline{sur} le segment) contenant $\delta I$.

\uline{NB} 1) Si $x,y\in L$, alors $x,y$ sont adjacents sss $\exists I\in S$ tel que $\{x,y\}=\delta I$. Donc, la donnée (4) \deleted{est inutile} résulte des autres (dans 7) il faut dire que si $\delta I=\{x,y\}$, $x$ et $y$ sont disjoints\ldots)

2) Je présume que si $x,y\in L$ sont adjacents, il y aura au plus deux segments $I$ tels que $\delta I=\{x,y\}$. (il y en aura effectivement deux sss $x,y$ appartiennent à une ``composante connexe'' de $\mathcal X$ qui est un (modèle d'un) cercle.

3) Je ne sais s'il faut supposer que $I$ est \unsure{connue} quand on connaît $\delta I$ et $\tilde I$, i.e. si on peut \deleted{définir par \unsure{le|un}} \inserted{$S\subset\unsure{\mathcal P}_2(\unsure{L_0|L})\times\unsure{\mathcal P}(L)$}

\uline{Axiomes} \inserted{\unsure{supplém.}} (à voir)

\deleted{On dit que la forme est de dim 0}

\margin{\deleted{\unsure{*}}}
$\unsure{E_1|E_2}$ Dans $L_0$, la relation ``$x$ et $y$ ne sont \uline{pas} disjoints'' est une relation d'équivalence.

L'ensemble quotient est appelé l'ens des \uline{points} \inserted{critiques, ou de branchement} de la forme.
\deleted{(\unsure{de})}
(\deleted{\unsure{NB}} plusieurs lieux critiques \unsure{différents} peuvent définir le même point critique)
}
\rotatebox{90}{\reflectbox{\unsure{EXTRA STRONG}}}

\rotatebox{90}{\reflectbox{\unsure{EXTRA STRONG}}}

\mspage{2}
